\chapter{Monte Carlo Simulation}

\section{Monte Carlo Event Generation}

Monte Carlo (MC) event generators provide a probabilistic description of proton--proton collisions, translating perturbative QCD calculations into fully simulated final states that can be processed through detector simulation. 
Modern generators combine fixed-order matrix-element calculations with parton showers, non-perturbative hadronisation models and phenomenological descriptions of additional soft interactions. 
The generation of an event proceeds through a sequence of well-defined stages, starting from the hard partonic interaction and ending with stable particles entering the detector simulation.

\subsection*{Hard Scattering and QCD Factorisation}

In hadron collisions, the fundamental interaction occurs between the constituent partons of the incoming protons. 
The calculation of production cross sections relies on the QCD factorisation theorem, which separates the short-distance perturbative dynamics from the long-distance structure of the proton. 
The inclusive cross section for producing a final state $X$ can be written as

\begin{equation}
\sigma_{pp\to X}
=
\sum_{i,j}
\int dx_1 dx_2\,
f_i(x_1,\mu_F)\,
f_j(x_2,\mu_F)\,
\hat{\sigma}_{ij\to X}(x_1,x_2,\mu_R,\mu_F),
\end{equation}

where $f_i(x,\mu_F)$ are the parton distribution functions (PDFs), describing the probability of finding a parton of type $i$ carrying a momentum fraction $x$ of the proton at the factorisation scale $\mu_F$. 
The partonic cross section $\hat{\sigma}_{ij\to X}$ is calculable in perturbative QCD and depends on both the renormalisation scale $\mu_R$ and the factorisation scale $\mu_F$. 
The renormalisation scale controls the value of the strong coupling constant $\alpha_s(\mu_R)$, while the factorisation scale defines the separation between perturbative and non-perturbative contributions.

In practice, the partonic cross section is expanded in powers of $\alpha_s$. 
Leading-order (LO) calculations provide a qualitative description of the process, while next-to-leading-order (NLO) calculations significantly improve the accuracy of total cross sections and reduce the dependence on the choice of scales. 
Residual scale variations are commonly used to estimate theoretical uncertainties.

Matrix elements at LO or NLO accuracy are computed using dedicated generators such as \textsc{MadGraph5\_aMC@NLO}, \textsc{Powheg}, or \textsc{Sherpa}. 
These tools evaluate the hard scattering amplitudes and generate parton-level configurations corresponding to the Born and higher-order contributions.

\subsection*{Parton Showers}

Fixed-order calculations do not account for multiple soft and collinear emissions, which can lead to large logarithmic corrections in certain regions of phase space. 
Parton shower algorithms approximate the all-orders resummation of leading logarithmic contributions by iteratively generating successive $1 \to 2$ branchings. 
These emissions are ordered in an evolution variable, such as transverse momentum or emission angle, and are described probabilistically through Sudakov form factors.

Parton showers simulate radiation from both the initial state (initial-state radiation, ISR) and the final state (final-state radiation, FSR). 
They are implemented in general-purpose event generators such as \textsc{Pythia8}, which employs a transverse-momentum-ordered shower, \textsc{Herwig7}, which uses angular ordering to preserve QCD coherence, and the dipole-based shower implemented in \textsc{Sherpa}. 
The parton shower stage provides a realistic description of additional jet activity beyond the fixed-order matrix-element calculation.

\subsection*{Matching and Merging}

To consistently combine fixed-order matrix elements with parton showers, matching and merging procedures are required. 
Without such prescriptions, radiation generated by the shower could double count contributions already included in the matrix-element calculation.

Matching schemes such as the \textsc{Powheg} and \textsc{MC@NLO} methods ensure NLO accuracy for inclusive observables while avoiding double counting of emissions. 
In the \textsc{Powheg} approach, the hardest emission is generated first using the exact matrix element, while softer emissions are subsequently produced by the parton shower. 
The \textsc{MC@NLO} method employs a subtraction formalism to combine NLO calculations with the shower.

For processes involving multiple hard jets, multi-leg merging schemes are used to combine matrix elements with different parton multiplicities. 
Algorithms such as CKKW or MEPS@NLO, as implemented in \textsc{Sherpa}, provide an improved description of additional hard radiation and are particularly important for high-multiplicity final states.

\subsection*{Hadronisation and Multiple Parton Interactions}

Following the perturbative evolution, coloured partons must be transformed into colourless hadrons due to confinement. 
This non-perturbative transition is modelled phenomenologically. 
The Lund string model, used in \textsc{Pythia8}, describes colour fields as strings that fragment into hadrons when stretched. 
The cluster model, used in \textsc{Herwig7}, groups colour-connected partons into colour-singlet clusters that subsequently decay into hadrons.

Additional soft activity in the event arises from multiple parton interactions (MPI), where several parton--parton scatterings occur within the same proton--proton collision. 
MPI and underlying-event activity are described using dedicated phenomenological models and tuned parameters, such as the A14 tune in \textsc{Pythia8}.

\subsection*{Heavy-Particle Decays and Event Finalisation}

Unstable particles produced during the perturbative or hadronisation stages are subsequently decayed according to their branching fractions. 
Spin correlations in heavy-particle decays can be preserved using dedicated tools such as \textsc{MadSpin}, while heavy-flavour hadron decays are often handled by programs such as \textsc{EvtGen}. 

The final output of the event generator consists of stable particles, which serve as input to the detector simulation. 
The modelling choices at each stage of the generation procedure introduce systematic uncertainties that are evaluated through variations of scales, PDF sets, shower models and tunes, as discussed in the following sections.
